\documentclass[tikz,border=10pt]{standalone}
\usetikzlibrary{positioning, arrows.meta, calc, babel}
\usepackage[utf8]{inputenc}
\usepackage[T1]{fontenc}
\usepackage{lmodern}
\usepackage[spanish,es-tabla]{babel}
\usepackage{amsmath, amssymb, bm}

\begin{document}
\begin{tikzpicture}[
  scale=2.5,
  align=center,
  font=\sffamily\small,
  % Estilos de ejes
  axisB/.style={-{Stealth[scale=1.2]}, line width=1.5pt},
  % Estilos del dron
  arm/.style={line width=4.5pt, color=gray!70},
  rotor/.style={fill=gray!20, draw=black!70, line width=1.2pt},
  propeller/.style={fill=white!90!black, opacity=0.4, draw=black!30},
  % Estilo de flechas de giro (un solo arco de ~270° por rotor)
  spinCW/.style={-{Stealth[scale=1.0]}, line width=1.2pt, color=red!70!black},
  spinCCW/.style={-{Stealth[scale=1.0]}, line width=1.2pt, color=blue!70!black},
]

  % Coordenadas de los rotores (en metros de cuerpo, escalados para la figura)
  % En cuerpo FRD: x_B hacia delante (arriba), y_B hacia la derecha (derecha)
  \coordinate (R0) at (0.75, 0.75);   % Delantero-Derecho
  \coordinate (R1) at (-0.75, 0.75);  % Delantero-Izquierdo
  \coordinate (R2) at (0.75, -0.75);  % Trasero-Derecho
  \coordinate (R3) at (-0.75, -0.75); % Trasero-Izquierdo
  \coordinate (CM) at (0, 0);
  \coordinate (zAxis) at (-1.35, 0);

  % ==========================================
  % EJES DE REFERENCIA CUERPO B (FRD)
  % ==========================================
  % x_B apunta hacia delante (arriba)
  \draw[axisB, color=cyan!80!black] (CM) -- (0, 1.3) node[above, font=\sffamily\bfseries] {$x_B$ (Delante)};
  % y_B apunta a la derecha (derecha)
  \draw[axisB, color=magenta!80!black] (CM) -- (1.3, 0) node[right, font=\sffamily\bfseries] {$y_B$ (Derecha)};
  % z_B apunta hacia abajo (hacia dentro de la página) — callout a la izquierda
  \node[draw=violet!80!black, circle, line width=1.2pt, inner sep=2pt] (zSym) at (zAxis) {};
  \draw[violet!80!black, line width=1.2pt] ($(zSym)+(-0.08,0.08)$) -- ($(zSym)+(0.08,-0.08)$);
  \draw[violet!80!black, line width=1.2pt] ($(zSym)+(-0.08,-0.08)$) -- ($(zSym)+(0.08,0.08)$);
  \node[below=0.12cm of zSym, font=\sffamily\bfseries, color=violet!80!black] {$z_B$ (Abajo)};

  % ==========================================
  % BRAZOS DEL CUADRICÓPTERO (Configuración en X)
  % ==========================================
  \draw[arm] (R1) -- (R2);
  \draw[arm] (R0) -- (R3);

  % Centro de masas
  \draw[fill=white, draw=black, line width=1.5pt] (CM) circle (0.15cm);
  \fill[black] (0,0) -- (0,0.15) arc (90:180:0.15) -- (0,0) -- (-0.15,0) arc (180:270:0.15) -- (0,0);
  \node[anchor=south west, font=\sffamily\scriptsize\bfseries] at (0.05, 0.05) {CM};

  % ==========================================
  % ROTORES, HÉLICES Y GIROS
  % ==========================================

  % --- ROTOR 0: Delantero-Derecho ---
  \draw[rotor] (R0) circle (0.15cm);
  \draw[propeller] (R0) ellipse (0.35cm and 0.1cm);
  \node[draw=black!60, fill=white, circle, inner sep=1pt, font=\sffamily\scriptsize] at (R0) {0};
  \node[anchor=west, xshift=1.0cm, font=\sffamily\scriptsize, text=black!70] at (R0)
    {$\bm{r}_0 = (0.17, 0.17)$ m};
  \draw[spinCW, ->] ($(R0) + (0, 0.28)$) arc (90:-180:0.28);

  % --- ROTOR 1: Delantero-Izquierdo ---
  \draw[rotor] (R1) circle (0.15cm);
  \draw[propeller] (R1) ellipse (0.35cm and 0.1cm);
  \node[draw=black!60, fill=white, circle, inner sep=1pt, font=\sffamily\scriptsize] at (R1) {1};
  \node[anchor=east, xshift=-1.0cm, font=\sffamily\scriptsize, text=black!70] at (R1)
    {$\bm{r}_1 = (0.17, -0.17)$ m};
  \draw[spinCCW, ->] ($(R1) + (0, 0.28)$) arc (90:360:0.28);

  % --- ROTOR 2: Trasero-Derecho ---
  \draw[rotor] (R2) circle (0.15cm);
  \draw[propeller] (R2) ellipse (0.35cm and 0.1cm);
  \node[draw=black!60, fill=white, circle, inner sep=1pt, font=\sffamily\scriptsize] at (R2) {2};
  \node[anchor=west, xshift=1.0cm, font=\sffamily\scriptsize, text=black!70] at (R2)
    {$\bm{r}_2 = (-0.17, 0.17)$ m};
  \draw[spinCCW, ->] ($(R2) + (0, 0.28)$) arc (90:360:0.28);

  % --- ROTOR 3: Trasero-Izquierdo ---
  \draw[rotor] (R3) circle (0.15cm);
  \draw[propeller] (R3) ellipse (0.35cm and 0.1cm);
  \node[draw=black!60, fill=white, circle, inner sep=1pt, font=\sffamily\scriptsize] at (R3) {3};
  \node[anchor=east, xshift=-1.0cm, font=\sffamily\scriptsize, text=black!70] at (R3)
    {$\bm{r}_3 = (-0.17, -0.17)$ m};
  \draw[spinCW, ->] ($(R3) + (0, 0.28)$) arc (90:-180:0.28);

  % ==========================================
  % LEYENDA EXPLICATIVA (única referencia a s_i)
  % ==========================================
  \node[draw=gray!30, fill=gray!3, rounded corners=2pt, below=3.0cm of CM, inner sep=5pt, minimum width=2.4cm, font=\sffamily\scriptsize] (leyenda) {
    \textbf{Signo del par de reacción $s_i$ (en guiñada):} \\
    \textcolor{red!70!black}{$\bm{\curvearrowright}$} Horario (CW), $s_i = -1$\\
    \textcolor{blue!70!black}{$\bm{\curvearrowleft}$} Antihorario (CCW), $s_i = +1$
  };

\end{tikzpicture}
\end{document}